\chapter{Conclusiones y recomendaciones}\label{cap:conclusiones}

\section{Respuesta a la pregunta}

Entre avisos de alquiler de departamentos en 13 distritos de Lima con dirección exacta, y bajo el supuesto de que área, distribución, antigüedad y ubicación capturan lo que determina a la vez la cochera y el precio, \textbf{tener cochera incluida eleva el precio de oferta en 7.5\,\% (IC 95\,\%: 3.7--11.4\,\%), unos S/\,199 al mes}. La prima es casi nula donde la cochera es habitual y llega al 18\,\% donde escasea. El resultado es estable ante cambios de estimador y de definición del tratamiento, y sensible a los errores de especificación que advierte el DAG. Su fragilidad es considerable: un confusor tan influyente como la antigüedad bastaría para llevarlo cerca de cero.

\section{Conclusiones por objetivo}

\begin{table}[H]
\centering\small
\caption{Objetivo, evidencia y conclusión verificable}\label{tab:correspondencia}
\begin{tabularx}{\textwidth}{>{\raggedright\arraybackslash}p{2.4cm}>{\raggedright\arraybackslash}p{3.2cm}X}
\toprule
Objetivo & Evidencia & Conclusión \\
\midrule
OE1 Auditoría & Tabla~\ref{tab:exclusiones} & 2\,220 avisos válidos de 3\,739; 243 duplicados y 30\,\% de direcciones imprecisas detectados \\
OE2 Identificación & Figura~\ref{fig:dag}, tabla~\ref{tab:roles}, figura~\ref{fig:soporte} & Estimando defendible en el soporte común (95\,\% de la muestra, grupos balanceados); amenaza: calidad del edificio \\
OE3 Estimación & Tabla~\ref{tab:principal_c} & +7.5\,\% frente a 63\,\% bruto; los estimadores coinciden. \textbf{H1 se confirma} \\
OE4 Variación y sensibilidad & Tabla~\ref{tab:especificaciones_c}, figuras~\ref{fig:sensibilidad_c} y \ref{fig:gate_c} & Los controles importan más que el estimador (\textbf{H2}); efecto heterogéneo por zona (\textbf{H3}); frágil frente a confusores (\textbf{H4}) \\
OE5 Medición con IA & Tabla~\ref{tab:extraccion} & Medición barata (0.76 USD) y trazable; exploratoria hasta validarla \\
OE6 Validación & Tabla~\ref{tab:sim_b} & Los estimadores funcionan sin confusión omitida; la cota binaria es optimista \\
OE7 Trazabilidad & Anexo~\ref{anx:inventario} & Cada resultado es rastreable a su notebook; falta un comando único \\
\bottomrule
\end{tabularx}
\end{table}

\paragraph{Tres tipos de afirmaciones.}
\begin{itemize}[itemsep=1pt]
\item \textbf{Observacional:} los avisos con cochera cuestan un 63\,\% más, pero casi todo es tamaño y zona.
\item \textbf{Causal condicionada a los supuestos:} la cochera incluida eleva el alquiler de oferta entre 3.7 y 11.4\,\%, y más donde escasea.
\item \textbf{Hipótesis futuras:} que la heterogeneidad responda a la escasez y a la normativa, y que la prima de oferta se traslade al contrato.
\end{itemize}

\section{Límites del trabajo}\label{sec:limites}

\begin{itemize}[leftmargin=*,itemsep=1pt]
\item \textbf{Identificación:} la ignorabilidad solo puede acotarse, y la calidad del edificio no se observa.
\item \textbf{Medición:} la IA no tiene validación humana. Además, ``sin mención'' se trata como ``sin cochera'', y el 72\,\% de las descripciones podría estar truncado.
\item \textbf{Tratamiento:} es sí/no; las cocheras múltiples (26\,\% de los tratados) se agrupan y las de costo extra se excluyen.
\item \textbf{Sensibilidad:} la cota binaria es optimista.
\item \textbf{Población:} oferta, no contrato; un portal y un mes; solo direcciones exactas, que excluyen el 30\,\% de los avisos.
\end{itemize}

\section{Recomendaciones}

\begin{enumerate}[itemsep=1pt]
\item \textbf{Validar la medición} con los 249 avisos etiquetados a mano y corregir el efecto con el estimador de \citet{egami2023}.
\item \textbf{Medir la calidad del edificio} desde el texto (ascensor, acabados, seguridad) para reducir la amenaza principal.
\item \textbf{Ampliar los datos} a más meses y portales, y \textbf{consolidar el repositorio} con un comando único reproducible.
\end{enumerate}
